% 本文件由 LaTeX 排版 Agent 根据有效 translation.json 自动生成。
% author=Ignatius of Antioch; work_id=0105; title=依纳爵致玛格尼西亚人书

\chapter{依纳爵致玛格尼西亚人书}

% structure: p0001 source=h1 target=omitted reason=page-title-already-rendered-by-parent

% structure: p0002 source=h2 target=section reason=relative-heading-rank; base=h2

\section{问候}

依纳爵，又名德奥佛鲁斯，致在玛格尼西亚（靠近米安得河）的教会，蒙天主圣父的恩宠，在耶稣基督我们的救主内，我问候这教会，并愿她在天主圣父和耶稣基督内充满幸福。\par

% structure: p0004 source=h2 target=section reason=relative-heading-rank; base=h2

\section{第一章　写信的缘由}

听闻你们那井然有序的虔诚之爱，我大为喜乐，并决意在耶稣基督的信德中与你们交谈。因为作为一位被视作堪当承受那最尊贵名号的人，在我所佩带的这些锁链中，我举荐众教会，并在其中祈求耶稣基督（我们生命的永恒泉源）及信德与爱德（无物能与之相比）的肉身与精神的合一，尤其是耶稣与圣父的合一；在祂内，如果我们忍受了此世王子的所有攻击并逃脱它们，我们将享见天主。\par

% structure: p0006 source=h2 target=section reason=relative-heading-rank; base=h2

\section{第二章　我为你们的使者而喜乐}

既然我已因着你们最可敬的主教达玛斯，以及你们可敬的长老巴苏斯和阿颇罗尼，还有我的同仆、执事索提奥（我愿永远享受他的友谊，因为他顺从主教如同顺从天主的恩宠，顺从长老团如同顺从耶稣基督的法律）而得以见到你们，[现在我写信给你们]。\par

% structure: p0008 source=h2 target=section reason=relative-heading-rank; base=h2

\section{第三章　尊敬你们年轻的主教}

你们也不应因他的年轻而过分亲近地对待你们的主教，反倒要向他致以所有的敬意，顾念天主圣父的权能；正如我所知道的可敬的长老们所做的，他们不因[他们的主教]显见的年轻外表而草率论断，反而因自己在天主内谨慎，顺从于他，或者更确切地说，不是顺从他，而是顺从耶稣基督的父，我们众人的主教。因此，你们应当以不虚伪的方式顺从[你们的主教]，以尊敬那愿我们[如此行]的那一位，因为那不顺从的人并非[藉此行为]欺骗那可见的主教，而是试图嘲弄那不可见者。所有此类行为所指向的并非人，而是那位知晓一切隐秘的天主。\par

% structure: p0010 source=h2 target=section reason=relative-heading-rank; base=h2

\section{第四章　有些人不义地离主教而行}

因此，不单要被称为基督徒，更要实实在在是基督徒，正如有些人确然称某人为主教，却离了他而行一切事。在我看来，这样的人似乎没有良好的良心，因为他们没有按照诫命坚定地聚集在一起。\par

% structure: p0012 source=h2 target=section reason=relative-heading-rank; base=h2

\section{第五章　如此之人的命运皆是死亡}

既然万物皆有终结，这两样就同时摆在我们面前——死亡与生命；并且各人将往自己的地方去。因为正如有两种钱币，一种属天主，另一种属世界，每种钱币都印有其特殊的印记，[此处亦然。]不信的人属于这个世界；但相信的人，在爱德中拥有天主圣父的印记，藉着耶稣基督，若不准备好死于祂的苦难，祂的生命就不在我们内。\par

% structure: p0014 source=h2 target=section reason=relative-heading-rank; base=h2

\section{第六章　保持和谐}

因此，既然我在前面所提到的人中，看见了你们整个团体在信德与爱德中，我劝勉你们努力以属神的和谐行事，以你们的主教代表天主之位，以你们的长老代表宗徒团之位，连同你们的执事——他们最为我所亲爱，并被托付了耶稣基督的职务，祂在时间开始之前就与圣父同在，并在末期被启示出来。你们所有人，效法这同一属神的行为，彼此尊重，不让任何人按肉身看待邻人，而要在耶稣基督内彼此恒久相爱。你们中间不可有任何分裂你们的事物，却要与你们的主教以及那些管理你们的人联合，作为你们不朽的模版与明证。\par

% structure: p0016 source=h2 target=section reason=relative-heading-rank; base=h2

\section{第七章　不可离了主教与长老行事}

因此，正如主没有圣父就无所作为，与祂合一，既不凭自己也不藉宗徒，同样，你们也不可离了主教与长老行事。你们也不要企图独自认为某事合理适切；反而要聚集在同一处，有同一祈祷、同一恳求、同一心意、同一望德，在无玷污的爱德与喜乐中。只有一位耶稣基督，无物比祂更卓越。所以，你们当一同跑进唯一的圣殿里，到唯一的祭台前，到唯一的耶稣基督面前，祂出自唯一的圣父，并与祂同在，且已归向唯一的圣父。\par

% structure: p0018 source=h2 target=section reason=relative-heading-rank; base=h2

\section{第八章　警惕虚假道理}

不要被怪异的道理所迷惑，也不要被无益的古老传说所欺。因为如果我们仍按犹太人的法律生活，那就是承认我们尚未领受恩宠。因为最神圣的先知们是按耶稣基督生活的。为此原因，他们也曾受迫害，被祂的恩宠所感动，以完全说服不信的人，使他们相信只有一位天主，祂已藉着祂的子耶稣基督显现了自己，耶稣基督是祂永恒的圣言，并非出自沉默，并在一切事上中悦了那派遣祂的一位。\par

% structure: p0020 source=h2 target=section reason=relative-heading-rank; base=h2

\section{第九章　让我们与基督同活}

因此，如果那些在古代秩序中长大的人已获得了新望德，不再守安息日，而生活在主日的遵守中，在此日我们的生命藉着祂和祂的死亡再次萌发——有些人否认这一点，我们却因这奥迹获得了信德，并因此忍受，好使我们得以成为耶稣基督（我们唯一的老师）的门徒——我们怎能离了祂而活呢？连先知们自己，在圣神中，也曾以祂为老师而等待祂。因此，他们所正当地等待的那一位，既已来临，就使他们从死人中复活。\BibleRef{玛 27:52}\par

% structure: p0022 source=h2 target=section reason=relative-heading-rank; base=h2

\section{第十章　谨防犹太化}

因此，我们不可对祂的仁慈无动于衷。因为倘若祂按我们的行为酬报我们，我们就不复存在了。所以，既然成了祂的门徒，就让我们学习按照基督信仰的原则生活吧。因为凡以此外的任何名称被称呼的，都不是出于天主。因此，你们要丢弃那邪恶的、旧的、发酸的酵母，变成新的酵母，即耶稣基督。当在祂内被盐渍，免得你们中有人败坏，因为你们必因自己的味道而被定罪。宣认耶稣基督却度犹太化的生活，是荒谬的。因为基督信仰并非拥抱犹太教，而是犹太教拥抱基督信仰，好使一切相信的舌头都聚集到天主那里。\par

% structure: p0024 source=h2 target=section reason=relative-heading-rank; base=h2

\section{第十一章 我写这些事是为了告诫你们}

我所亲爱的，我说这些事，并非我知道你们中有任何人处于这种状态；而是，作为你们中最小的一个，我愿预先提醒你们，免得你们落在虚假教义的钩子上，而是要在有关诞生、苦难和复活的事上获得充分的把握，这些事发生在般雀·比拉多执政时期，由我们的希望耶稣基督真实而确定地完成，\BibleRef{弟前 1:1}愿你们中无人偏离这希望。\par

% structure: p0026 source=h2 target=section reason=relative-heading-rank; base=h2

\section{第十二章 你们胜过我}

但愿我在各方面都蒙受你们，倘若我果真堪当！因为我虽被捆绑，却不配与你们任何自由的人相比。我知道你们并不自高自大，因为你们有耶稣基督在你们内。而且，当我称赞你们时，我更加知道你们珍惜谦逊的精神；正如经上所记载：\Quote{义人是他自己的控告者。}\BibleRef{箴 18:17}\par

% structure: p0028 source=h2 target=section reason=relative-heading-rank; base=h2

\section{第十三章 要在信德与合一中立定}

因此，你们要努力在主和宗徒的教导中立定，好使你们所做的一切事，在肉身和灵魂、在信德和爱德、在圣子、圣父和圣神、在起初和终结中，都顺利亨通；与你们那位极其可敬的主教，以及你们司铎团那编制精巧的属灵冠冕，并那些按天主心意的执事一起。你们要服从主教，也彼此服从，如同耶稣基督按肉身服从父，宗徒们服从基督、父和圣神一样，好使肉身和灵魂都得以合一。\par

% structure: p0030 source=h2 target=section reason=relative-heading-rank; base=h2

\section{第十四章 请求你们的祈祷}

我知道你们充满天主，因此只是简略地劝勉你们。请在你们的祈祷中纪念我，使我得以到达天主那里；也纪念在叙利亚的教会，我连从其得名都不配，因为我需要你们在天主内的共同祈祷和爱，好使在叙利亚的教会堪当因你们的教会而获得安慰。\par

% structure: p0032 source=h2 target=section reason=relative-heading-rank; base=h2

\section{第十五章 致候}

从斯米纳来的厄弗所人——我也是从那里写信给你们——他们在这里为天主的荣耀，正如你们一样，他们在一切事上安慰了我，他们问候你们，连同斯米纳的主教波利卡普。其余的教会，为尊崇耶稣基督，也问候你们。愿你们在天主的和谐中平安，你们已获得了那不可分离的神，就是耶稣基督。\par

\SourceRecord{Translated by Alexander Roberts and James Donaldson.；Ante-Nicene Fathers；Vol. 1.；Edited by Alexander Roberts, James Donaldson, and A. Cleveland Coxe.；Buffalo, NY: Christian Literature Publishing Co.,；1885.；<http://www.newadvent.org/fathers/0105.htm>.}
